\documentclass[11pt]{article}

\usepackage[a4paper,margin=1in]{geometry}
\usepackage{amsmath,amssymb,mathtools}
\usepackage{booktabs}
\usepackage{enumitem}
\usepackage{hyperref}
\usepackage{microtype}
\usepackage{siunitx}

\hypersetup{
  colorlinks=true,
  linkcolor=blue,
  urlcolor=blue,
  citecolor=blue
}

\sisetup{
  per-mode=symbol,
  detect-all=true
}

\title{A Reduced Model for Bragg-Interferometer Phase Noise Induced by Raman Velocity-Selection Noise}
\author{Project note for the local Python analysis tool}
\date{}

\begin{document}
\maketitle

\begin{abstract}
We document a reduced calculation model for estimating how shot-to-shot fluctuations of the Raman-selected velocity class propagate into the phase of a first-order Bragg Mach--Zehnder atom interferometer. The implementation targets the current ${}^{87}\mathrm{Rb}$ configuration with a $\pi/2$--$\pi$--$\pi/2$ pulse sequence, Gaussian Bragg pulses, and a Gaussian velocity-selected ensemble. Two observables are emphasized: the total interferometer fringe phase and an operational pulse-only Bragg diffraction phase obtained by suppressing dark-time phase accumulation while keeping the same finite-duration pulses. The document also defines the exact quantities shown in the four panels of the desktop interface: selected velocity distributions, phase-versus-center curves, phase-noise-versus-Raman-noise curves, and contrast-versus-center curves.
\end{abstract}

\section{Introduction and Scope}
The physical question is the following: after Raman velocity selection, each experimental shot prepares a narrow but fluctuating velocity class. Even if the microscopic origin of the fluctuation is not modeled, the resulting shot-to-shot variation of the selected mean velocity changes the effective Doppler detuning seen by the Bragg pulses. Because finite-duration Bragg pulses have detuning-dependent transfer amplitudes and phases, the mean velocity fluctuation is converted into a fluctuation of the interferometer output phase.

The present implementation is intentionally restricted to a controlled first model:
\begin{itemize}[leftmargin=1.5em]
  \item atomic species: ${}^{87}\mathrm{Rb}$,
  \item interferometer sequence: $\pi/2$--$\pi$--$\pi/2$,
  \item Bragg order: first order,
  \item Bragg pulses: Gaussian envelopes in an effective two-level model,
  \item selected atoms: one-dimensional Gaussian detuning distribution,
  \item shot-to-shot Raman noise: Gaussian noise on the center detuning of that selected distribution,
  \item cloud size, spatial inhomogeneity, and position--velocity correlations: neglected.
\end{itemize}

This scope is sufficient to quantify, within a simplified but internally consistent model, how a Raman center-noise RMS maps to a Bragg-induced phase-noise RMS.

\section{Default Experimental Configuration}
The default user-interface values correspond to the current experimental discussion:
\begin{center}
\begin{tabular}{ll}
\toprule
Quantity & Default value \\
\midrule
Atomic species & ${}^{87}\mathrm{Rb}$ \\
Bragg order & $n=1$ \\
Pulse sequence & $\pi/2$--$\pi$--$\pi/2$ \\
Beam-splitter duration & \SI{25}{\micro\second} \\
Mirror duration & \SI{50}{\micro\second} \\
Bragg pulse shape & Gaussian \\
Dark time & \SI{8}{\milli\second} \\
Selected-distribution FWHM & \SI{12}{\kilo\hertz} \\
Reference Raman center-noise RMS & \SI{4}{\kilo\hertz} \\
Nominal center detuning & \SI{0}{\kilo\hertz} relative to Bragg resonance \\
\bottomrule
\end{tabular}
\end{center}

The statement that the Raman-selected $+\hbar k$ class corresponds to approximately \SI{15}{\kilo\hertz} in the Raman configuration is treated here as calibration context. In the current Bragg solver, the dynamical variable is the detuning relative to the nominal Bragg resonance of the selected class.

\section{Effective Bragg Pulse Model}
\subsection{Two-level Hamiltonian}
For a given Doppler-equivalent two-photon detuning $\delta$ and pulse phase $\varphi$, the first-order Bragg pulse is modeled in the basis
\[
\bigl\{\, |0\rangle,\; |2\hbar k\rangle \,\bigr\}
\]
by the effective Hamiltonian
\begin{equation}
\frac{H(t;\delta,\varphi)}{\hbar}
=
\frac{1}{2}
\begin{pmatrix}
-\delta & \Omega(t)e^{-i\varphi} \\
\Omega(t)e^{+i\varphi} & \delta
\end{pmatrix}.
\label{eq:effective_hamiltonian}
\end{equation}
For the first two pulses the code uses $\varphi=0$, while the final pulse phase is the scanned interferometer phase.

\subsection{Gaussian pulse envelope and area calibration}
The Bragg envelope is written as
\begin{equation}
\Omega(t)=\Omega_0\exp\!\left(-\frac{t^2}{2\tau^2}\right),
\label{eq:gaussian_envelope}
\end{equation}
where $\tau$ is the RMS temporal width used internally by the solver. Since the UI may specify the pulse duration under different conventions, the code first converts the user entry to $\tau$.

The pulse is truncated numerically at $|t|\le N_\sigma\tau$. The peak Rabi frequency $\Omega_0$ is set by the target pulse area,
\begin{equation}
\theta
=
\int_{-N_\sigma \tau}^{+N_\sigma \tau}\Omega(t)\,dt
=
\Omega_0\tau\sqrt{2\pi}\,\mathrm{erf}\!\left(\frac{N_\sigma}{\sqrt{2}}\right),
\label{eq:pulse_area}
\end{equation}
with $\theta=\pi/2$ for the beam splitter and $\theta=\pi$ for the mirror.

\subsection{Time-discretized exact propagator}
For each short time step $\Delta t$ the Hamiltonian is approximated as constant, and the corresponding propagator is evaluated analytically. Writing
\begin{equation}
M(\delta,\Omega)=
\begin{pmatrix}
-\delta & \Omega \\
\Omega & \delta
\end{pmatrix},
\end{equation}
with phase absorbed into the chosen basis for the relevant step, the step propagator is
\begin{equation}
U_{\mathrm{step}}(\delta,\Omega,\Delta t)
=
\cos\!\left(\frac{\Omega_{\mathrm{eff}}\Delta t}{2}\right)I
-
i\,\frac{\sin\!\left(\frac{\Omega_{\mathrm{eff}}\Delta t}{2}\right)}{\Omega_{\mathrm{eff}}}
M(\delta,\Omega),
\label{eq:step_propagator}
\end{equation}
where
\begin{equation}
\Omega_{\mathrm{eff}}=\sqrt{\delta^2+\Omega^2}.
\end{equation}
The unitary of a finite Gaussian pulse is obtained by ordered multiplication of Eq.~\eqref{eq:step_propagator} along the pulse.

\section{Free Evolution and Pulse Sequence}
In the same rotating-frame description, dark-time evolution is represented by
\begin{equation}
U_{\mathrm{free}}(T)
=
\begin{pmatrix}
e^{+i\delta T/2} & 0 \\
0 & e^{-i\delta T/2}
\end{pmatrix}.
\label{eq:free_evolution}
\end{equation}
For a fixed detuning class $\Delta$, the full three-pulse sequence is therefore
\begin{equation}
U_{\mathrm{MZ}}(\Delta,\phi)
=
U_{\mathrm{BS}}(\Delta,\phi)
U_{\mathrm{free}}(T)
U_{\mathrm{M}}(\Delta)
U_{\mathrm{free}}(T)
U_{\mathrm{BS}}(\Delta),
\label{eq:full_sequence}
\end{equation}
where $\phi$ is the phase of the final beam splitter.

\section{Raman-Selected Ensemble and Shot-to-Shot Noise}
The Raman-selected atoms are represented by a Gaussian detuning distribution
\begin{equation}
g(\Delta;\mu,\sigma)
=
\frac{1}{\sqrt{2\pi}\sigma}
\exp\!\left[-\frac{(\Delta-\mu)^2}{2\sigma^2}\right],
\label{eq:gaussian_distribution}
\end{equation}
where $\mu$ is the shot-dependent center detuning. The selected-distribution width is parameterized by the experimental FWHM,
\begin{equation}
\sigma = \frac{\mathrm{FWHM}}{2\sqrt{2\ln 2}}.
\label{eq:fwhm_to_sigma}
\end{equation}

The Raman center noise is then modeled as
\begin{equation}
\mu_{\mathrm{shot}}\sim\mathcal{N}\!\left(\mu_0,\sigma_{\mathrm{noise}}^2\right),
\label{eq:center_noise}
\end{equation}
with nominal center $\mu_0$ and shot-to-shot RMS $\sigma_{\mathrm{noise}}$.

For ensemble observables, averages over Eq.~\eqref{eq:gaussian_distribution} are evaluated numerically by Gauss--Hermite quadrature.

\section{Fringe Coefficient, Phase, and Contrast}
\subsection{Single-detuning fringe coefficient}
Let the input state be $|\psi_{\mathrm{in}}\rangle=|0\rangle$. For a fixed detuning class $\Delta$, the first beam splitter, the mirror, and the two dark-time propagators produce a state immediately before the final beam splitter,
\begin{equation}
|\psi_{\mathrm{pre}}(\Delta)\rangle
=
a(\Delta)|0\rangle+b(\Delta)|2\hbar k\rangle.
\label{eq:pre_final_state}
\end{equation}
The last beam splitter is then the element that converts the two incoming amplitudes $a(\Delta)$ and $b(\Delta)$ into a measurable interference fringe.

Under the phase convention used in the implementation, the $|0\rangle$ output amplitude after the final beam splitter can be written as
\begin{equation}
\mathcal{A}_0(\phi\mid\Delta)
=
\alpha(\Delta)+\beta(\Delta)e^{-i\phi},
\label{eq:output_amplitude_decomposition}
\end{equation}
where
\begin{equation}
\alpha(\Delta)=U_{\mathrm{BS},00}(\Delta)\,a(\Delta),
\qquad
\beta(\Delta)=U_{\mathrm{BS},01}(\Delta)\,b(\Delta).
\label{eq:alpha_beta_def}
\end{equation}
Equation~\eqref{eq:output_amplitude_decomposition} is precisely the statement that the final beam splitter mixes the two interferometer arms into the detected output port.

The output probability is therefore
\begin{align}
P_0(\phi\mid\Delta)
&=
\left|\mathcal{A}_0(\phi\mid\Delta)\right|^2 \\
&=
\left|\alpha(\Delta)\right|^2
+
\left|\beta(\Delta)\right|^2
+
\alpha(\Delta)\beta^*(\Delta)e^{i\phi}
+
\alpha^*(\Delta)\beta(\Delta)e^{-i\phi}.
\label{eq:expanded_output_probability}
\end{align}
This expression can always be recast as
\begin{equation}
P_0(\phi\mid\Delta)=A(\Delta)+2\,\mathrm{Re}\!\left[c(\Delta)e^{i\phi}\right],
\label{eq:single_detuning_fringe}
\end{equation}
with
\begin{equation}
A(\Delta)=\left|\alpha(\Delta)\right|^2+\left|\beta(\Delta)\right|^2,
\qquad
c(\Delta)=\alpha(\Delta)\beta^*(\Delta).
\label{eq:single_c_definition}
\end{equation}
Hence $c(\Delta)$ is not an additional phenomenological parameter; it is the complex cross term generated when the two pre-final amplitudes interfere at the last beam splitter. Its argument determines the fringe shift for that single detuning class,
\begin{equation}
\phi_{\mathrm{tot}}(\Delta)=-\arg\!\left[c(\Delta)\right].
\label{eq:single_class_phase}
\end{equation}

In the implementation, this construction appears explicitly. The code first computes $|\psi_{\mathrm{pre}}(\Delta)\rangle$, then forms the two partial output amplitudes of Eq.~\eqref{eq:alpha_beta_def}, and finally evaluates $c(\Delta)=\alpha(\Delta)\beta^*(\Delta)$.

\subsection{Ensemble-averaged total phase}
The observed interferometer output is not the single-class quantity $c(\Delta)$, but its average over the selected ensemble. The ensemble-averaged fringe coefficient is
\begin{equation}
\bar{c}_{\mathrm{tot}}(\mu)
=
\int g(\Delta;\mu,\sigma)\,c_{\mathrm{tot}}(\Delta)\,d\Delta,
\label{eq:ensemble_total_c}
\end{equation}
and the corresponding total interferometer phase is
\begin{equation}
\Phi_{\mathrm{tot}}(\mu)
=
-\arg\!\left[\bar{c}_{\mathrm{tot}}(\mu)\right].
\label{eq:ensemble_total_phase}
\end{equation}
The associated ensemble contrast is defined by
\begin{equation}
C_{\mathrm{tot}}(\mu)
=
\frac{2\left|\bar{c}_{\mathrm{tot}}(\mu)\right|}{\bar{A}_{\mathrm{tot}}(\mu)},
\qquad
\bar{A}_{\mathrm{tot}}(\mu)=\int g(\Delta;\mu,\sigma)A_{\mathrm{tot}}(\Delta)\,d\Delta.
\label{eq:total_contrast}
\end{equation}
The key point is that the measured phase comes from the argument of the \emph{ensemble-averaged} complex fringe coefficient $\bar c_{\mathrm{tot}}(\mu)$, not from averaging the phases $\phi_{\mathrm{tot}}(\Delta)$ directly.

\subsection{Operational pulse-only diffraction phase}
To isolate the phase generated by the finite-duration Bragg pulses themselves, the tool also defines a pulse-only sequence obtained by setting the free-evolution operator to the identity while keeping the same pulse unitaries. The pulse-only pre-final state is
\begin{equation}
|\psi_{\mathrm{pre}}^{(\mathrm{diff})}\rangle
=
U_{\mathrm{M}}(\Delta)
U_{\mathrm{BS}}(\Delta)
|0\rangle.
\label{eq:pulse_only_state}
\end{equation}
Repeating the same fringe-coefficient extraction gives
\begin{equation}
\bar{c}_{\mathrm{diff}}(\mu)
=
\int g(\Delta;\mu,\sigma)\,c_{\mathrm{diff}}(\Delta)\,d\Delta,
\label{eq:ensemble_diff_c}
\end{equation}
\begin{equation}
\Phi_{\mathrm{diff}}(\mu)
=
-\arg\!\left[\bar{c}_{\mathrm{diff}}(\mu)\right],
\label{eq:ensemble_diff_phase}
\end{equation}
\begin{equation}
C_{\mathrm{diff}}(\mu)
=
\frac{2\left|\bar{c}_{\mathrm{diff}}(\mu)\right|}{\bar{A}_{\mathrm{diff}}(\mu)}.
\label{eq:diff_contrast}
\end{equation}
This quantity is model dependent, but it is well-defined within the present solver and directly measures the detuning-dependent phase shift produced by the Bragg pulses in the absence of dark-time phase accumulation.

\section{Propagation of Raman Center Noise into Phase Noise}
The tool reports two mappings from Raman center-noise RMS to phase-noise RMS.

\subsection{Linearized propagation}
At the nominal center $\mu_0$, the small-noise estimate is
\begin{equation}
\sigma_{\phi,\mathrm{tot}}^{(\mathrm{lin})}
\approx
\left|\frac{d\Phi_{\mathrm{tot}}}{d\mu}\right|_{\mu=\mu_0}
\sigma_{\mathrm{noise}},
\label{eq:linear_total_noise}
\end{equation}
\begin{equation}
\sigma_{\phi,\mathrm{diff}}^{(\mathrm{lin})}
\approx
\left|\frac{d\Phi_{\mathrm{diff}}}{d\mu}\right|_{\mu=\mu_0}
\sigma_{\mathrm{noise}}.
\label{eq:linear_diff_noise}
\end{equation}
In the code, these slopes are computed numerically by a centered finite difference in $\mu$.

\subsection{Monte Carlo propagation}
The non-linear estimate is obtained by sampling
\begin{equation}
\mu_j\sim\mathcal{N}(\mu_0,\sigma_{\mathrm{noise}}^2),
\qquad j=1,\dots,N_{\mathrm{MC}},
\end{equation}
and evaluating the wrapped phase excursions relative to the nominal shot,
\begin{equation}
\delta\Phi_{\mathrm{tot},j}
=
\arg\!\left\{\exp\!\left[i\left(\Phi_{\mathrm{tot}}(\mu_j)-\Phi_{\mathrm{tot}}(\mu_0)\right)\right]\right\},
\label{eq:wrapped_total_excursion}
\end{equation}
\begin{equation}
\delta\Phi_{\mathrm{diff},j}
=
\arg\!\left\{\exp\!\left[i\left(\Phi_{\mathrm{diff}}(\mu_j)-\Phi_{\mathrm{diff}}(\mu_0)\right)\right]\right\}.
\label{eq:wrapped_diff_excursion}
\end{equation}
The corresponding RMS values are
\begin{equation}
\sigma_{\phi,\mathrm{tot}}^{(\mathrm{MC})}
=
\operatorname{Std}\!\left[\delta\Phi_{\mathrm{tot},j}\right],
\qquad
\sigma_{\phi,\mathrm{diff}}^{(\mathrm{MC})}
=
\operatorname{Std}\!\left[\delta\Phi_{\mathrm{diff},j}\right].
\label{eq:mc_noise_definitions}
\end{equation}
This makes the role of $c(\Delta)$ in the Monte Carlo chain completely explicit:
\begin{equation}
c(\Delta)
\longrightarrow
\bar c(\mu_j)
\longrightarrow
\Phi(\mu_j)=-\arg\!\left[\bar c(\mu_j)\right]
\longrightarrow
\sigma_\phi.
\label{eq:c_to_sigma_chain}
\end{equation}
For each synthetic shot center $\mu_j$, the code first averages the single-class interference kernel $c(\Delta)$ over the shifted Gaussian ensemble, then extracts the shot phase from the argument of the resulting complex number, and finally computes the standard deviation across many such shots.

\section{Interpretation of the Four Plots Produced by the UI}
The graphical interface displays four panels. Their mathematical definitions are summarized here so that the figures can be discussed in a paper-like way.

\subsection{Panel 1: selected velocity distribution}
The upper-left panel shows the nominal and shifted selected distributions. The nominal profile is
\begin{equation}
I_{\mathrm{nom}}(\Delta)
\propto
\exp\!\left[-\frac{(\Delta-\mu_0)^2}{2\sigma^2}\right],
\label{eq:nominal_distribution_plot}
\end{equation}
while the reference shifted profile shown by the UI corresponds to a $+1$ RMS center shift,
\begin{equation}
I_{+1\sigma}(\Delta)
\propto
\exp\!\left[-\frac{(\Delta-\mu_0-\sigma_{\mathrm{noise}})^2}{2\sigma^2}\right].
\label{eq:shifted_distribution_plot}
\end{equation}
These curves are not themselves interferometer observables; they visualize the Raman-selected input ensemble and the meaning of the shot-to-shot center fluctuation.

\subsection{Panel 2: phase change versus shot center detuning}
The upper-right panel plots the phase change relative to the nominal center. For the total interferometer phase, the plotted quantity is
\begin{equation}
\Delta\Phi_{\mathrm{tot}}(\mu)
=
\operatorname{unwrap}\!\left[\Phi_{\mathrm{tot}}(\mu)\right]
-
\operatorname{unwrap}\!\left[\Phi_{\mathrm{tot}}(\mu_0)\right].
\label{eq:delta_total_phase_plot}
\end{equation}
Likewise, for the pulse-only diffraction phase,
\begin{equation}
\Delta\Phi_{\mathrm{diff}}(\mu)
=
\operatorname{unwrap}\!\left[\Phi_{\mathrm{diff}}(\mu)\right]
-
\operatorname{unwrap}\!\left[\Phi_{\mathrm{diff}}(\mu_0)\right].
\label{eq:delta_diff_phase_plot}
\end{equation}
This panel directly shows how sensitive the interferometer phase is to shot-to-shot shifts of the selected mean velocity.

\subsection{Panel 3: phase noise versus Raman center-noise amplitude}
The lower-left panel maps the Raman center-noise RMS $\sigma_{\mathrm{noise}}$ to the resulting phase-noise RMS. Four curves are shown:
\begin{align}
\sigma_{\phi,\mathrm{tot}}^{(\mathrm{lin})}(\sigma_{\mathrm{noise}})
&=
\left|\frac{d\Phi_{\mathrm{tot}}}{d\mu}\right|_{\mu_0}\sigma_{\mathrm{noise}},
\label{eq:plot_total_linear}
\\
\sigma_{\phi,\mathrm{tot}}^{(\mathrm{MC})}(\sigma_{\mathrm{noise}})
&=
\operatorname{Std}\!\left[\delta\Phi_{\mathrm{tot},j}\right],
\label{eq:plot_total_mc}
\\
\sigma_{\phi,\mathrm{diff}}^{(\mathrm{lin})}(\sigma_{\mathrm{noise}})
&=
\left|\frac{d\Phi_{\mathrm{diff}}}{d\mu}\right|_{\mu_0}\sigma_{\mathrm{noise}},
\label{eq:plot_diff_linear}
\\
\sigma_{\phi,\mathrm{diff}}^{(\mathrm{MC})}(\sigma_{\mathrm{noise}})
&=
\operatorname{Std}\!\left[\delta\Phi_{\mathrm{diff},j}\right].
\label{eq:plot_diff_mc}
\end{align}
Agreement between linearized and Monte Carlo curves indicates a locally linear response around $\mu_0$; deviation indicates a non-linear phase response over the noise scale being probed.

\subsection{Panel 4: contrast versus shot center detuning}
The lower-right panel plots the ensemble-averaged contrast versus the shot center detuning,
\begin{equation}
C_{\mathrm{tot}}(\mu)
=
\frac{2|\bar{c}_{\mathrm{tot}}(\mu)|}{\bar{A}_{\mathrm{tot}}(\mu)},
\qquad
C_{\mathrm{diff}}(\mu)
=
\frac{2|\bar{c}_{\mathrm{diff}}(\mu)|}{\bar{A}_{\mathrm{diff}}(\mu)}.
\label{eq:contrast_plot_definitions}
\end{equation}
This panel is important because phase sensitivity and contrast are coupled: a steep phase response accompanied by strong contrast loss may indicate that the selected distribution is moving into a spectrally unfavorable region of the Bragg response.

\section{Discussion of What the Four Plots Mean Physically}
Taken together, the four panels provide a compact diagnostic chain:
\begin{enumerate}[leftmargin=1.5em]
  \item Panel 1 defines the input ensemble and the magnitude of the imposed shot-to-shot Raman fluctuation.
  \item Panel 2 converts that input fluctuation into a phase-response curve in detuning space.
  \item Panel 3 turns the local or non-linear phase response into a directly useful noise metric, namely $\sigma_\phi$.
  \item Panel 4 tracks whether the same detuning excursion also degrades contrast, which is often a signature of operating away from optimal Bragg resonance.
\end{enumerate}
In a paper-style discussion, Panel 2 is the response function, Panel 3 is the propagated noise metric, and Panel 4 is the associated visibility penalty.

\section{Assumptions and Limitations}
The model is rigorous only within the reduced description above. It does \emph{not} yet include:
\begin{itemize}[leftmargin=1.5em]
  \item coupling to higher Bragg momentum orders,
  \item explicit multi-state diffraction phases,
  \item spatial intensity inhomogeneity and wavefront effects,
  \item initial cloud size or position distribution,
  \item Raman pulse dynamics or a microscopic model for the origin of the center noise,
  \item correlations between center fluctuations and width fluctuations.
\end{itemize}
Accordingly, the current results should be interpreted as a controlled first estimate of how a noisy Raman-selected mean velocity maps into total fringe phase noise and pulse-only diffraction phase noise.

\section{Code Organization}
The implementation used by the desktop application is organized as follows:
\begin{itemize}[leftmargin=1.5em]
  \item \texttt{main.py}: launches the desktop UI,
  \item \texttt{bragg\_phase\_noise/config.py}: model configuration,
  \item \texttt{bragg\_phase\_noise/physics.py}: Gaussian pulses and finite-step propagators,
  \item \texttt{bragg\_phase\_noise/analysis.py}: ensemble averages, phases, contrasts, and noise propagation,
  \item \texttt{bragg\_phase\_noise/ui.py}: Tkinter + Matplotlib interface.
\end{itemize}

\end{document}
